\guidechapter{The smallest transform step, applied on many axes}{Read before simultaneous butterfly layers.}

A two-point normalized Fourier transform sends
\[
 H_0(u,v)=\left(\frac{u+v}{2},\frac{u-v}{2}\right).
\]
The first output records an average; the second records a difference.
For $(u,v)=(6,2)$ the outputs are $(4,2)$. Recover the original by
$u=4+2$, $v=4-2$. The factor $1/2$ matters: this is not the unnormalized
Hadamard matrix, nor the unitary matrix with factor $1/\sqrt2$.

\explfigure{06-butterfly}{A butterfly and a two-axis layer.}{
The lower example uses the tensor product $H_0\otimes H_0$ on four entries.
One axis gives $[6,2,2,0]$, and the other gives $[4,1,2,1]$.
The entries are examples, not intermediate values from the large network.}

If every input has complex magnitude at most $M$, each output also does:
$|u\pm v|/2\le(|u|+|v|)/2\le M$. Thus the completed operation is
contractive in the maximum norm. Applying it on several address bits means
pairing records that differ in a selected bit, for every setting of the other
bits. A naive implementation makes one whole-array pass per selected axis.
The proposed layer algorithm organizes those operations jointly.

\paragraph{Why the paper temporarily uses a different kernel.}
Define $X$ to exchange two entries, $\mathsf S=\operatorname{diag}(1,i)$,
$a=(1+i)/2$, $b=(1-i)/2$, and $C=aI+bX$. The useful identities are
\[
 C^2=X,\quad C^4=I,\qquad H_0=b\mathsf S C\mathsf S.
\]
The kernel $C$ interacts well with the binary phase frames used to price
residual directions. Once its tensor layer is complete, diagonal phases and
a known dyadic complex factor produce the desired $H_0$ layer. These are
classical linear identities; the use of phase matrices does not require a
quantum computer.

\paragraph{One final rounding, not one rounding per internal gate.}
The internal values are exact Gaussian dyadics, numbers whose real and
imaginary parts are integers divided by powers of two. Internal scratch can
grow even though the completed $H_0$ layer is contractive. The proof therefore
needs guard bits for internal magnitudes and denominators. It executes the
whole layer exactly and truncates only at its boundary. Contraction alone
would not justify discarding bits after every scratch update.

\paragraph{Why selected bits create exceptional addresses.}
An efficient packed address update may agree with a desired binary basis
change on most addresses while carries disturb a small, explicitly described
set. The paper isolates that set and repairs it deterministically. The proof
needs agreement outside it, invariance of the exceptional set, and a bound on
the cost of extracting, sorting, and reinserting full records. ``Rare'' does
not mean random or ignorable: every input must still be correct.

\takeaway{The second saving batches arithmetic across axes. It must survive
basis changes, intermediate growth, and exact correction of exceptional records.}
\sourceanchor{05-layers.tex, eq:phase-C; eq:phase-basic-identities}
